\section{A finite Gaussian-fiber comparison}\label{sec:r-fiber}

We now bound the information retained after a block of exact Gaussian
measurements by comparing two ways to average the state likelihood along
a measurement fiber.  The comparison measure is finite but need not have
mass one.  Its marginal is an inverse-distance kernel involving the rows
actually used by the transition.  A common local-mass quantity then
connects that marginal to a projection independent of the new state.

We retain the measure and entropy conventions of
Section~\ref{sec:r-experiments}.  Constants denoted by \(K\) are
positive and absolute and may increase.  The local row counts are
\begin{equation}\label{rfib:parameters}
 n=\lfloor d/10\rfloor,\qquad
 k=\lfloor d/3\rfloor,\qquad
 \ell=\lfloor d/10\rfloor.
\end{equation}
Here \(n\) is the number of rows read in one block, \(k\) is the number
of rows in an auxiliary projection, and \(\ell\) is an integer moment
order.  We take \(d\) sufficiently large that \(\ell\ge2\) and
\(k-\ell-n\ge d/10\).

\subsection{Projection fibers and conditional information}

We first record the geometric facts used both in the block estimate and
at the terminal state.  If \(\sigma_m\) is uniform probability on
\(S^{m-1}\), then for \(m\ge2\), \(z\in S^{m-1}\), and \(0<r\le1/2\),
\begin{equation}\label{rfib:cap}
 \sigma_m(B(z,r))\le r^{m-1}.
\end{equation}
This is the centered-ball estimate of Lemma~\ref{lem:r-caps} in
dimension \(m\).

For a full-row-rank \(b\times d\) matrix \(C\), where \(1\le b\le d-2\),
put
\[
 P_C=C^{\mathsf T}(CC^{\mathsf T})^{-1}C,\qquad
 x_C(z)=C^{\mathsf T}(CC^{\mathsf T})^{-1}z .
\]
Write \(S(H)=\{u\in H:\|u\|=1\}\) for the unit sphere of a subspace
\(H\).  The conditional law of \(S\sim\sigma\) given \(CS=z\), denoted by
\(\sigma_{C,z}\), fixes \(P_CS=x_C(z)\) and is uniform on the remaining
sphere
\[
 x_C(z)+\sqrt{1-\|x_C(z)\|^2}\,
       S(\ker C),\qquad \|x_C(z)\|<1.
\]
This formula specifies a Borel probability kernel on the indicated
set.  For example, the uniform direction in \(\ker C\) can be realized
as \((I-P_C)Z/\|(I-P_C)Z\|\) with \(Z\) a standard Gaussian vector;
the map is Borel away from a null set.  Set the kernel to a fixed point mass off
the full-rank interior set.  To verify the formula, write a uniform
sphere point as a standard Gaussian vector divided by its length.
The squared lengths in the row space and its orthogonal complement are
independent chi-squared variables with \(b\) and \(d-b\) degrees of
freedom before division.  The first divided by their sum has the beta distribution with
parameters \(b/2\) and \((d-b)/2\), and the two directions are
independent and uniform.  Consequently the row-space projection, in
orthonormal coordinates, has density
\begin{equation}\label{rfib:projection-density}
 c_{d,b}(1-\|x\|^2)^{(d-b-2)/2}\quad(\|x\|<1),
 \qquad
 c_{d,b}=\frac{\Gamma(d/2)}
               {\pi^{b/2}\Gamma((d-b)/2)} .
\end{equation}
The density is zero outside the unit ball.  In particular, the conditional statement holds for almost every
observation.  Since
\(\|x_C(z)\|^2=z^{\mathsf T}(CC^{\mathsf T})^{-1}z\), the
Lebesgue density \(p_C\) of \(CS\) is
\begin{equation}\label{rfib:label-density}
 p_C(z)=
 \begin{cases}
 \displaystyle
 \frac{c_{d,b}}{\det(CC^{\mathsf T})^{1/2}}
 \bigl(1-\|x_C(z)\|^2\bigr)^{(d-b-2)/2},
       &\|x_C(z)\|<1,\\
 0,    &\text{otherwise}.
 \end{cases}
\end{equation}
Boundary values do not affect any of the integrals below.

For \(b=k\), the exponent in \eqref{rfib:projection-density} is
nonnegative.  Integrating its normalization over the ball of radius
\(1/2\) shows that its supremum times the volume of the unit \(k\)-ball
is at most \(2^k(4/3)^{d/2}\).  Thus, for every ball \(D\) of radius
\(r\) in the row space,
\begin{equation}\label{rfib:projection-ball}
 \Pr\{P_CS\in D\}\le 2^k(4/3)^{d/2}r^k\le e^{Kd}r^k.
\end{equation}
The estimate is for the unweighted uniform measure \(\sigma\).

Let \(V\) take values in a finite set \(\mathcal V\) and be jointly
distributed with \(S\).  Choose Borel versions of the state likelihoods
\[
 f_v(s)=\Pr\{V=v\mid S=s\},\qquad
 \bar f_v(C,z)=\int f_v(t)\,d\sigma_{C,z}(t).
\]
Modify the first family on a Borel \(\sigma\)-null set, if needed, so
that \(f_v\ge0\) and \(\sum_v f_v=1\) everywhere.
These are likelihoods, not probability densities normalized separately
for each state.  When \(C\) has \(k\) independent standard Gaussian rows
and is independent of \((S,V)\), disintegration identifies
\(\bar f_v(C,z)\) with \(\Pr\{V=v\mid C,CS=z\}\).  The conditional
relative-entropy chain rule therefore gives
\begin{equation}\label{rfib:potential}
 \begin{aligned}
 \Phi(V)&:=I(S;V\mid C,CS)
 =\mathbb E\log\frac{f_V(S)}{\bar f_V(C,CS)}\\
 &=H(V\mid C,CS)-H(V\mid S).
 \end{aligned}
\end{equation}
Here \(H\) denotes Shannon entropy in nats.  For the standard-Borel form of this chain rule, see
\cite[Section~3.3]{Austin2020}.  Nonnegativity of conditional relative
entropy and the displayed difference give
\(0\le\Phi(V)\le\log|\mathcal V|\).  Terms on
zero-probability state or fiber events are omitted.

\begin{proposition}[Information through a Gaussian block]\label{rfib:block}
Let \(S\sim\sigma\), let \(V\) be finite-valued, and let \(A\) have
\(n\) independent standard Gaussian rows, independently of \((S,V)\).
Suppose \(W\) takes at most \(N\) values, where \(N\ge1\) is an integer, and is obtained from
\((V,A,AS)\) by a Borel transition kernel.  If \(\log N\le d^2\), then
for the dimensions in \eqref{rfib:parameters} and all sufficiently
large \(d\),
\begin{equation}\label{rfib:drift}
 \Phi(W)\le\Phi(V)+Kd .
\end{equation}
The auxiliary matrix in each occurrence of \(\Phi\) is independent of
the pair consisting of the signal and the state in that occurrence.
\end{proposition}

The proof first compares the real law to a finite measure.  We next
identify the latter's signal/state marginal by a two-point Borel measure
identity.  Finally, two estimates with different geometric proofs compare
that marginal and the fresh-projection likelihood to one local mass.

\subsection{The finite measure and its exact marginal}

Write the transition probabilities as \(h_{vw}(A,AS)\).  Define
\begin{equation}\label{rfib:routing}
 g_w(A,t)=\sum_v f_v(t)h_{vw}(A,At),\qquad
 f_w(t)=\mathbb E_A g_w(A,t).
\end{equation}
Thus \(0\le g_w\le1\), \(0\le f_w\le1\), and the sums over \(w\) of
each of these functions are one.  The function \(f_w\) is precisely
the new state's likelihood given \(S=t\).  In particular neither
\(g_w\) nor \(f_w\) is divided by the marginal probability of state
\(w\).

Augment \(A\) by \(k-n\) independent Gaussian rows and call the
result \(C\), with \(A\) as its first \(n\) rows.  Relative to the
product measure \(\mathfrak m\) of the Gaussian law of \(C\),
\(\sigma(ds)\), and counting measure on the two state sets, the real law \(P\) and the comparison
measure \(Q\) on \((C,S,V,W)\) have densities
\begin{align}
 \frac{dP}{d\mathfrak m}(C,s,v,w)
   &=f_v(s)h_{vw}(A,As),\notag\\
 \frac{dQ}{d\mathfrak m}(C,s,v,w)
   &=p_C(Cs)\bar f_v(C,Cs)h_{vw}(A,As).
   \label{rfib:PQ}
\end{align}
The factor \(p_C(Cs)\) turns the conditional fiber average
\(\bar f_v(C,Cs)\) into the density at \(Cs\) of the projection of
\(f_v\,d\sigma\).  This unnormalized equal-label measure will yield the
inverse-distance kernel in Lemma~\ref{rfib:two-point}.
The routing weights in \(Q\) remain the actual weights
\(h_{vw}(A,As)\).  We have \(P\ll Q\): \(p_C(Cs)>0\) almost surely,
and a fiber on which \(\bar f_v(C,Cs)=0\) has \(f_v=0\) almost
everywhere on that fiber.

The comparison measure has strictly positive finite mass.  Indeed,
summing the routing probabilities and then the old-state likelihoods
gives
\[
 m_Q:=Q(\Omega)
 =\mathbb E_C\int p_C(Cs)\,d\sigma(s)
 =\mathbb E_C\int p_C(z)^2\,dz
 \le\mathbb E_C\sup_z p_C(z)<\infty.
\]
The last expectation is finite by \eqref{rfib:label-density}.
Gram--Schmidt on the Gaussian rows expresses
\(\det(CC^{\mathsf T})^{1/2}\) as a product of successive
perpendicular lengths with conditional chi laws of degrees
\(d,d-1,\ldots,d-k+1\).  A chi variable of degree \(\nu>1\) has
finite inverse first moment
\(2^{-1/2}\Gamma((\nu-1)/2)/\Gamma(\nu/2)\), and all these degrees
are greater than one.  Strict positivity follows from the positivity
of \(p_C\) on the interior of its support.

Use the signed finite-reference entropy
\(D_{\mathrm{fin}}(P\|Q)=\mathbb E_P\log(dP/dQ)\) from
Lemma~\ref{lem:r-finite-reference}.  Cancellation in
\eqref{rfib:PQ}, together with the independence of \(C\) from
\((S,V)\), yields
\begin{equation}\label{rfib:finite-entropy}
 D_{\mathrm{fin}}(P\|Q)=\Phi(V)+\mathbb E[-\log p_C(CS)]
 \le\Phi(V)+\frac{k}{2}\log(2\pi e).
\end{equation}
Here \(\log p_C(CS)\) is integrable: the logarithm of the determinant
has finite log moments by the same chi factorization, and the residual
factor in \eqref{rfib:label-density} has finite log moments by the beta
law.  For the inequality, let \(\varphi_k\) be the standard Gaussian
density.  Nonnegativity of \(D(p_C\,dz\|\varphi_k\,dz)\), integrated
over \(C\), gives
\(\mathbb E[-\log p_C(CS)]\le\mathbb E[-\log\varphi_k(CS)]\).
Unconditionally \(CS\) is standard Gaussian, because \(S\) has unit
length and is independent of the Gaussian rows.  This assertion is
not a conditional Gaussian assertion given \(C\) or the state.

It is necessary to normalize \(Q\) before invoking probability data
processing.  Set \(Q_0=Q/m_Q\), and let \(\pi\) be any measurable
map.  Both \(Q\) and its pushforward have mass \(m_Q\), so
\begin{align}
 D_{\mathrm{fin}}(P\|Q)&=D(P\|Q_0)-\log m_Q, \notag\\
 D_{\mathrm{fin}}(\pi_\#P\|\pi_\#Q)
   &=D(\pi_\#P\|\pi_\#Q_0)-\log m_Q \notag\\
   &\le D(P\|Q_0)-\log m_Q
     =D_{\mathrm{fin}}(P\|Q).                         \label{rfib:finite-dp}
\end{align}
No nonnegativity of \(D_{\mathrm{fin}}(P\|Q)\) is used.

Let \(q_w(s)\) be the density of the \((S,W)\) marginal of \(Q\)
relative to \(\sigma\) and counting measure.  The corresponding real
density is \(f_w(s)\).  Equations \eqref{rfib:finite-entropy} and
\eqref{rfib:finite-dp} imply
\begin{equation}\label{rfib:marginal-entropy}
 \mathbb E\log\frac{f_W(S)}{q_W(S)}
 \le\Phi(V)+Kd ,
\end{equation}
where this expectation, and all state expectations below, are under
the real law.  The densities \(q_w\) need not sum to a probability
density.

We next calculate this marginal in a form suitable for geometric
estimates.  Let \(\gamma_b\) denote the law of \(b\) unrestricted
standard Gaussian rows.  For a unit vector \(u\), let
\(\gamma_{b,\perp u}\) be the law of \(b\) independent rows with
Gaussian covariance \(I-uu^{\mathsf T}\), supported on \(u^\perp\).

\begin{lemma}[Two-point Gaussian fiber measure]\label{rfib:two-point}
Assume \(1\le k<d-1\).  For every bounded Borel function
\(H(C,s,t)\) on
\(\mathbb R^{k\times d}\times S^{d-1}\times S^{d-1}\),
\begin{align}
 &\int \gamma_k(dC)\int\sigma(ds)\,p_C(Cs)
       \int\sigma_{C,Cs}(dt)\,H(C,s,t) \notag\\
 &\quad=(2\pi)^{-k/2}\int\sigma(ds)\int\sigma(dt)\,
       \|s-t\|^{-k}
       \int\gamma_{k,\perp\widehat{s-t}}(dC)\,H(C,s,t),
                                                        \label{rfib:two-point-id}
\end{align}
where \(\widehat{s-t}=(s-t)/\|s-t\|\) off the diagonal.  The diagonal
has product measure zero, and its assigned value is immaterial.
Both sides define finite Borel measures.
\end{lemma}

\begin{proof}
For \(\delta>0\), let
\(\kappa_\delta(z)=(2\pi\delta^2)^{-k/2}
 e^{-\|z\|^2/(2\delta^2)}\).  First let \(H\) be bounded and
continuous, and integrate
\[
 H(C,s,t)\kappa_\delta(C(t-s))
 \quad\text{against}\quad
 \gamma_k(dC)\sigma(ds)\sigma(dt).
\]
Fix full-rank \(C\) and \(s\) with \(\|P_Cs\|<1\).  Disintegration
in \(z=Ct\) makes the integral in \(t\)
\[
 \int\kappa_\delta(z-Cs)\,p_C(z)
       \left(\int H(C,s,t)\,d\sigma_{C,z}(t)\right)dz.
\]
The function multiplied by \(\kappa_\delta\) is continuous at \(Cs\).
Indeed the displayed fiber parametrization is continuous at interior
observations, and the test is bounded and continuous.  It is bounded
in absolute value by \(\|H\|_\infty\sup_zp_C(z)\).  The Gaussian
approximate identity therefore converges to the left integrand in
\eqref{rfib:two-point-id}; the preceding inverse determinant estimate
makes this bound integrable in \(C\).

Alternatively fix \(s\ne t\), and integrate first in \(C\).
Writing \(v=t-s\), each row integral has total weight
\[
 (2\pi)^{-1/2}(\|v\|^2+\delta^2)^{-1/2}.
\]
After division by that weight, the row is a centered Gaussian with
covariance \(I-vv^{\mathsf T}/(\|v\|^2+\delta^2)\).  As
\(\delta\downarrow0\) this law converges to the Gaussian law on
\(v^\perp\).  The \(k\) rows remain independent, so the limit is the
right integrand in \eqref{rfib:two-point-id}.  Its absolute value is
dominated by
\((2\pi)^{-k/2}\|H\|_\infty\|s-t\|^{-k}\).
This is integrable under \(\sigma\otimes\sigma\): dyadic annuli about
\(s\), with \eqref{rfib:cap} for the small radii, give a convergent
geometric series because \(k<d-1\).

The two orders compute the same limit.  They prove equality on bounded
continuous tests of two finite Borel measures on the indicated Polish
space.  Uniqueness of finite Borel measures then proves the equality
for bounded Borel tests as well.  This last step is what permits
measurable routing functions in the application.
\end{proof}

On the fiber \(Ct=Cs\), the first \(n\) rows give \(At=As\).
Consequently, using \eqref{rfib:routing} without normalizing its
weights,
\[
 q_w(s)=\mathbb E_C\left[
   p_C(Cs)\int g_w(A,t)\,d\sigma_{C,Cs}(t)\right].
\]
Apply Lemma~\ref{rfib:two-point} to
\(H(C,s,t)=\psi(s)g_w(A,t)\) for bounded Borel \(\psi\).
The extra \(k-n\) rows do not enter \(g_w\), so integrating them out
leaves precisely the marginal law of the \(n\) actual rows.  We obtain,
for \(\sigma\)-almost every \(s\),
\begin{equation}\label{rfib:kernel}
 q_w(s)=(2\pi)^{-k/2}\int
   \|s-t\|^{-k}
   \mathbb E_{A\sim\gamma_{n,\perp\widehat{s-t}}}g_w(A,t)
   \,d\sigma(t).
\end{equation}
This is an equality of marginal densities of the finite measure \(Q\);
there is no additional normalization on its right-hand side.

\subsection{Two estimates for a common local mass}

Put \(a_w=\int f_w\,d\sigma=\Pr\{W=w\}\), and omit states with
\(a_w=0\).  Define
\begin{equation}\label{rfib:local-mass}
 L_{w,r}(s)=r^{-k}\int_{B(s,r)}f_w(t)\,d\sigma(t),
 \qquad L_w(s)=\sup_{r>0}L_{w,r}(s).
\end{equation}
For each \(r\) the integral is Borel in \(s\).  Rational radii suffice
in the supremum: approximate any radius from above and use continuity
from above of the finite measure \(f_w\sigma\).  Thus \(L_w\) is
Borel.  The diameter-two ball gives its lower bound, while
\eqref{rfib:cap} for \(r\le1/2\) and total mass at most one for larger
radii give the upper bound
\begin{equation}\label{rfib:local-bounds}
 2^{-k}a_w\le L_w(s)\le2^k .
\end{equation}
Choose the state-count-dependent cutoff
\begin{equation}\label{rfib:cutoff}
 J=1+\left\lceil\frac{\log_2N}{d-1-k}\right\rceil,\qquad
 r_*=2^{-J}.
\end{equation}
Then \(r_*\le1/2\), \(Nr_*^{d-1-k}\le1\), and \(J=O(d)\), since
\(d-1-k\ge d/2\) for large \(d\) and \(\log N\le d^2\).

\begin{lemma}[Comparison with an independent projection]\label{rfib:first-ratio}
Let \(C'\) have \(k\) independent standard Gaussian rows and be
independent of \((S,W)\).  Then
\begin{equation}\label{rfib:ratio-fresh}
 \mathbb E\frac{L_W(S)}{\bar f_W(C',C'S)}\le e^{Kd}.
\end{equation}
\end{lemma}

\begin{proof}
Fix a full-rank \(C'\) and a radius \(r\le2\).  Choose a finite
maximal \(r\)-separated set \(\{z_i\}\) on the sphere, and partition
the sphere into Borel cells \(E_i\subseteq B(z_i,r)\), breaking ties
by their indices.  The enlarged balls \(B(z_i,2r)\) overlap at most
\(5^d\) times: balls of radius \(r/2\) about their separated centers
have disjoint interiors, and all those meeting a given point in an
enlarged ball lie inside the ambient ball of radius \(5r/2\) about
that point.

On \(E_i\),
\[
 L_{w,r}(s)\le r^{-k}\int_{B(z_i,2r)}f_w\,d\sigma.
\]
Write \(R\) for the row space of \(C'\), \(p_R\) for the density in
\eqref{rfib:projection-density}, and \(\sigma_x\) for the corresponding
uniform fiber law.  Conditioning on \(P_{C'}S=x\) gives
\begin{align*}
 \int_{E_i}\frac{f_w(s)}{\bar f_w(C',C's)}\,d\sigma(s)
 &=\int_R p_R(x)
   \frac{\int_{E_i}f_w(s)\,d\sigma_x(s)}
        {\bar f_w(C',C'x)}\,dx\\
 &\le\Pr\{P_{C'}S\in B(P_{C'}z_i,r)\}
 \le e^{Kd}r^k .
\end{align*}
The quotient is defined to be zero on zero-denominator fibers; on all
other fibers its numerator is at most its denominator.  The numerator
vanishes on fibers outside the displayed projected ball.  The last
probability is under the unweighted prior and is bounded by
\eqref{rfib:projection-ball}.  Multiplying the last two estimates and
summing the cells, with their \(5^d\) overlap, yields
\begin{equation}\label{rfib:fixed-radius}
 \int\frac{f_w(s)L_{w,r}(s)}
           {\bar f_w(C',C's)}\,d\sigma(s)
 \le e^{Kd}a_w .
\end{equation}

For \(r\ge r_*\), one of the grid radii
\(2,1,1/2,\ldots,r_*\) gives
\(L_{w,r}\le2^k L_{w,\mathrm{grid}}\); radii greater than two are
already bounded by the radius-two expression.  For \(r<r_*\),
\eqref{rfib:cap} gives
\(L_{w,r}\le r^{d-1-k}\le r_*^{d-1-k}\).
A second conditioning on the projection gives
\(\int f_w/\bar f_w(C',C'\cdot)\,d\sigma\le1\).
We may therefore sum \eqref{rfib:fixed-radius} over the grid and over
states, obtaining
\[
 \sum_w\int\frac{f_wL_w}{\bar f_w(C',C'\cdot)}\,d\sigma
 \le 2^k(J+2)e^{Kd}\sum_w a_w+Nr_*^{d-1-k}
 \le e^{Kd}.
\]
The bound is uniform in full-rank \(C'\), and averaging proves the
claim.
\end{proof}

The other comparison is with the actual finite-measure marginal:
\begin{equation}\label{rfib:ratio-kernel}
 \mathbb E\frac{q_W(S)}{L_W(S)}\le e^{Kd}.
\end{equation}
We prove it now.  For each positive-mass state introduce the finite
measure
\[
 d\nu_w(s)=\frac{f_w(s)}{L_w(s)}\,d\sigma(s).
\]
Equation \eqref{rfib:local-bounds} gives
\(\nu_w(S^{d-1})\le2^k\).  More locally, for every ambient center
\(z\) and \(r>0\),
\begin{equation}\label{rfib:nu-ball}
 \nu_w(B(z,r))\le(2r)^k .
\end{equation}
Indeed, if \(m=\int_{B(z,r)}f_w\,d\sigma>0\), every \(s\) in that
ball on the sphere has \(L_w(s)\ge(2r)^{-k}m\); integration of
\(f_w/L_w\) over the ball gives the bound.  The case \(m=0\) is
immediate.

Inserting \eqref{rfib:kernel} shows that the left side of
\eqref{rfib:ratio-kernel} is \(\sum_w\int q_w\,d\nu_w\).  First take
the portion with \(\|s-t\|\le r_*\).  Uniformly in \(s\), dyadic
annuli and \eqref{rfib:cap} give
\begin{equation}\label{rfib:near}
 \int_{\{\|s-t\|\le r_*\}}\|s-t\|^{-k}\,d\sigma(t)
 \le\frac{2^k r_*^{d-1-k}}
          {1-2^{-(d-1-k)}}.
\end{equation}
Using \(g_w\le1\), summing the at most \(N\) masses
\(\nu_w(S^{d-1})\le2^k\), and retaining the prefactor in
\eqref{rfib:kernel}, this portion is at most
\[
 (2\pi)^{-k/2}\,
 \frac{N\,2^{2k} r_*^{d-1-k}}
      {1-2^{-(d-1-k)}}\le e^{Kd}.
\]
Thus the state count in the near singular part is paid by the cutoff
chosen in \eqref{rfib:cutoff}.

For the remaining distances use the annuli
\(r/2<\|s-t\|\le r\) with outer radii
\(r=2,1,\ldots,2r_*\).  Fix \(w,t,r\), and push
\(r^{-k}\nu_w\) restricted to that annulus to the unit direction
\(u=(s-t)/\|s-t\|\); call the resulting measure \(\mu\).
It has total mass at most \(2^k\) by \eqref{rfib:nu-ball}.
For a linear subspace \(E\) of dimension \(0\le j\le\ell-1\) and
\(0<\eta\le1\), its directions satisfy
\begin{equation}\label{rfib:tube}
 \mu\{u:\operatorname{dist}(u,E)\le\eta\}
 \le3^j4^k\eta^{k-j}.
\end{equation}
In fact the corresponding \(s\)'s lie within distance \(\eta r\) of
the affine subspace \(t+E\), with their projections there in the
radius-\(r\) ball about \(t\).  That \(j\)-ball has an
\(\eta r\)-net of at most \((3/\eta)^j\) points by ambient packing.
Balls of radius \(2\eta r\) about these points cover the \(s\)'s.
Applying \eqref{rfib:nu-ball} and multiplying by \(r^{-k}\) proves
\eqref{rfib:tube}.

\Needspace{6\baselineskip}
\begin{lemma}[Averaging perpendicular Gaussian laws]\label{rfib:mixture}
Let \(\mu\) be a finite measure on \(S^{d-1}\) of mass at most \(2^k\)
satisfying \eqref{rfib:tube} for every linear subspace of dimension
\(j\le\ell-1\).  Define the finite measure on \(n\times d\) matrices
\[
 \lambda=\int\gamma_{n,\perp u}\,d\mu(u).
\]
This is a Borel measure: a perpendicular matrix can be realized as
\(Z(I-uu^{\mathsf T})\) with \(Z\sim\gamma_n\).
Then every bounded Borel function \(H\) on this matrix space satisfies
\begin{equation}\label{rfib:mixture-dual}
 \left|\int H\,d\lambda\right|
 \le e^{Kd}\|H\|_{L^{\ell/(\ell-1)}(\gamma_n)} .
\end{equation}
\end{lemma}

\begin{proof}
For \(E\) of dimension \(j\le\ell-1\), layer-cake integration of
\eqref{rfib:tube} gives the inverse moment
\begin{equation}\label{rfib:inverse-distance}
 \int\operatorname{dist}(u,E)^{-n}\,d\mu(u)
 \le2^k+\frac{n\,3^j4^k}{k-j-n}\le e^{Kd}.
\end{equation}
Here \(\operatorname{dist}(u,E)\le1\); its layer-cake integral below
one is \(n\int_0^1\eta^{-n-1}
\mu\{\operatorname{dist}(u,E)\le\eta\}\,d\eta\).
The denominator is positive because
\(k-j-n\ge k-\ell+1-n\ge d/10\).

For \(u_1,\ldots,u_\ell\), let
\(G=(\langle u_i,u_j\rangle)_{i,j\le\ell}\).
The tube estimate, sent to \(\eta\downarrow0\), shows that each
proper span of at most \(\ell-1\) preceding directions has
\(\mu\)-mass zero.  Thus \(G\) is invertible for
\(\mu^{\otimes\ell}\)-almost every tuple.  Gram--Schmidt gives
\[
 \det G=\prod_{i=1}^{\ell}
 \operatorname{dist}\bigl(u_i,\operatorname{span}(u_1,\ldots,u_{i-1})
                    \bigr)^2 .
\]
Integrate \(u_\ell,u_{\ell-1},\ldots,u_1\) in that order and apply
\eqref{rfib:inverse-distance} at each step.  This proves
\begin{equation}\label{rfib:inverse-gram}
 \int(\det G)^{-n/2}\prod_{i=1}^{\ell}d\mu(u_i)
 \le e^{Kd\ell}.
\end{equation}

It remains to pass from this determinant moment to singular row laws.
For \(0<\delta\le1\), set
\[
 \begin{aligned}
 \alpha_\delta&=\Pr\{|Z|\le\delta\},\qquad Z\sim N(0,1),\\
 D_\delta(A)&=\int
 \frac{\mathbf 1_{\{\|Au\|_\infty\le\delta\}}}{\alpha_\delta^n}
 \,d\mu(u).
 \end{aligned}
\]
There is an absolute \(c>0\) such that \(\alpha_\delta\ge c\delta\).
Expanding the integer \(\ell\)-th moment under \(\gamma_n\) produces
one tuple \(u_1,\ldots,u_\ell\).  For an invertible Gram matrix \(G\),
the projections of each unrestricted Gaussian row onto those directions
have density at most
\((2\pi)^{-\ell/2}(\det G)^{-1/2}\).  The box of allowed projections
has volume \((2\delta)^\ell\).  Row independence and
\eqref{rfib:inverse-gram} therefore give
\[
 \mathbb E_{\gamma_n}D_\delta^\ell
 \le
 \left(\frac{(2\delta)^\ell(2\pi)^{-\ell/2}}
             {\alpha_\delta^\ell}\right)^n e^{Kd\ell}
 \le e^{Kd\ell}.
\]
The null set of dependent tuples was removed above.  Thus
\(\|D_\delta\|_{L^\ell(\gamma_n)}\le e^{Kd}\) uniformly in
\(\delta\).

The measure \(D_\delta\,d\gamma_n\) is the mixture, with mixing
measure \(\mu\), of the laws conditioned on
\(\|Au\|_\infty\le\delta\).  For each \(u\), this conditioning
preserves the independent perpendicular components and makes all
components along \(u\) tend to zero.  For every bounded continuous
test \(H\), bounded convergence in \(u\) consequently gives
\(\int H D_\delta\,d\gamma_n\longrightarrow\int H\,d\lambda\).
Hölder's inequality proves \eqref{rfib:mixture-dual} for those tests.
Both \(\gamma_n\) and \(\lambda\) are finite Euclidean Borel
measures.  Bounded continuous functions are dense in
\(L^{\ell/(\ell-1)}(\gamma_n+\lambda)\); approximation in this
space passes the inequality to every bounded Borel \(H\).
In particular no absolute continuity of the singular mixture
\(\lambda\) was assumed before the dual estimate.
\end{proof}

For the far-annulus contribution to \(\sum_w\int q_w\,d\nu_w\),
apply Lemma~\ref{rfib:mixture} to the annular measure above and
\(H(A)=g_w(A,t)\).  Since \(0\le g_w\le1\),
\[
 \int\mathbb E_{\gamma_{n,\perp u}}g_w(A,t)\,d\mu(u)
 \le e^{Kd}\bigl(\mathbb E_{\gamma_n}g_w(A,t)\bigr)^{1-1/\ell}
 =e^{Kd}f_w(t)^{1-1/\ell}.
\]
On its annulus, \(\|s-t\|^{-k}\le2^kr^{-k}\).
The definition of \(\mu\) thus bounds the annular integration in
\(s\) in \eqref{rfib:kernel} by \(2^ke^{Kd}f_w(t)^{1-1/\ell}\),
apart from the displayed kernel prefactor.  Concavity gives, pointwise
in \(t\),
\[
 \sum_w f_w(t)^{1-1/\ell}
 \le N^{1/\ell}\left(\sum_wf_w(t)\right)^{1-1/\ell}
 \le N^{1/\ell}\le e^{Kd}.
\]
There are \(J+1=O(d)\) far annuli.  Summing them and integrating in
\(t\), then adding the near contribution, proves
\eqref{rfib:ratio-kernel}.

\subsection{Logarithmic combination}

We can now finish Proposition~\ref{rfib:block}.  All logarithms in the
following split are integrable under the real \((S,W)\) law and the
independent \(C'\).  Indeed
\[
 \mathbb E[-\log f_W(S)]=H(W\mid S)\le\log N,\qquad
 \mathbb E[-\log\bar f_W(C',C'S)]=H(W\mid C',C'S)\le\log N.
\]
Also \eqref{rfib:local-bounds} and
\(\mathbb E[-\log a_W]=H(W)\le\log N\) make \(\log L_W(S)\)
integrable.  Finally \(P_{S,W}\ll Q_{S,W}\).  For any finite
measure \(Q'\) and probability \(P'\ll Q'\), the negative part of
its log density ratio has bound
\[
 \int_{\{r<1\}}r\log(1/r)\,dQ'\le Q'(\Omega)/e,
 \qquad r=\frac{dP'}{dQ'}.
\]
Applied to the marginal comparison, this and the finite upper bound
\eqref{rfib:marginal-entropy} make
\(\log(f_W(S)/q_W(S))\) integrable.  Hence \(\log q_W(S)\) is
integrable as well, and \(q_W(S)>0\) almost surely under the real law.

Using the fresh \(C'\) in the new-state potential, we may therefore
write
\begin{align*}
 \Phi(W)
 &=\mathbb E\log\frac{f_W(S)}{q_W(S)}
   +\mathbb E\log\frac{q_W(S)}{L_W(S)}\\
 &\quad+\mathbb E\log\frac{L_W(S)}{\bar f_W(C',C'S)}\\
 &\le \Phi(V)+Kd
   +\log\mathbb E\frac{q_W(S)}{L_W(S)}\\
 &\quad+\log\mathbb E\frac{L_W(S)}{\bar f_W(C',C'S)}
 \le\Phi(V)+Kd .
\end{align*}
The first inequality uses \eqref{rfib:marginal-entropy} and Jensen;
the last uses \eqref{rfib:ratio-kernel} and
Lemma~\ref{rfib:first-ratio}.  This proves the proposition.  The
extra \(k-n\) rows allowed both local comparisons to use exponent
\(k\), whereas only the \(n\) actual rows entered the perpendicular
mixture moment.


Subsection~\ref{subsec:r-fiber-application} applies the estimate to a
learner and combines it with stopping and terminal accuracy to obtain a
sample lower bound.
